\documentclass{beamer}

% Load style (configured via Makefile/scripts)
\usepackage{\jobname_style}

% Load template preamble (shared, not duplicated)
\include{../../00-preamble}

% Presentation metadata (example-specific, humorous)
% Auto-generated presentation metadata
% Edit config/presentation.conf and rerun ./configure-presentation.sh

\title{The ten most important results of astrology}
\author{Dr. Foo Bar}
\institute{Hogwarts University}
\date{\today}
\newcommand{\presentationevent}{International Conference on Nonsense}
\newcommand{\presentationlocation}{Moon Base Alpha, Lunar Surface}
\newcommand{\email}{foo@mail.com}


\begin{document}

% ------------------------------------------------------------------
% Title slide
% ------------------------------------------------------------------
\begin{frame}
    \titlepage
\end{frame}

% ------------------------------------------------------------------
% Introduction
% ------------------------------------------------------------------
\section{Introduction}

\begin{frame}{Motivation}
    \begin{itemize}
        \item Astrology has long resisted experimental verification
        \item Nevertheless, it continues to make strong predictions
        \item We present its \emph{ten most important results}
    \end{itemize}
\end{frame}

% ------------------------------------------------------------------
% Math / Physics examples
% ------------------------------------------------------------------
\section{Theoretical Framework}

\begin{frame}{Fundamental Assumption}
    The dynamics of celestial influence are governed by
    \[
        \mathcal{F} \propto \frac{\text{Mercury in retrograde}}{\text{Peer review}}
    \]
\end{frame}

\begin{frame}{Equation of Motion}
    Assuming natural units ($c = \hbar = 1$), the horoscope field
    $\phi_{\mathrm{astro}}$ satisfies
    \[
        \Box \phi_{\mathrm{astro}} + m_{\star}^2 \phi_{\mathrm{astro}}
        = \lambda\, \cos(\theta_{\mathrm{zodiac}})
    \]
\end{frame}

\begin{frame}{Dimensional Analysis}
    The prediction strength scales as
    \[
        P \sim \mathcal{O}(1)\times \left(\frac{\text{confidence}}{\text{evidence}}\right)^n,
        \qquad n \gg 1
    \]
\end{frame}

% ------------------------------------------------------------------
% Physics-style slide
% ------------------------------------------------------------------
\section{Results}

\begin{frame}{Experimental Validation}
    \begin{itemize}
        \item Dataset size: $N = 0$
        \item Statistical significance: $5\sigma$ (by convention)
        \item Systematic uncertainties: neglected
    \end{itemize}
\end{frame}

% ------------------------------------------------------------------
% Closing slide
% ------------------------------------------------------------------
\section{Conclusion}

\begin{frame}{Summary}
    \begin{itemize}
        \item Demonstrated mathematical typesetting
        \item Illustrated physics-style notation
        \item No actual astrology was harmed in this example
    \end{itemize}
\end{frame}

\end{document}
